\documentclass[12pt, twoside]{article}
\input{../preamble}

\fancyhead[LE]{\thepage}
\fancyhead[RO]{Markscheme}
\fancyhead[LO]{La Scuola d'Italia / Dr.\ Huson / IB Math AA SL \\* 13 May 2026}
\fancyfoot[C]{\thepage}

\begin{document}
\subsubsection*{11.18 Classwork: Series --- Markscheme}

\begin{enumerate}[itemsep=0.35cm]

%--- Problem 1: Geometric sums and sums of partial sums [5 marks] ---
%--- Source: AA SL 2024 May P1 TZ1 Q6 ---
\item[1.] [Maximum mark: 5]

\textbf{(a)} Find $S_n$. \hfill [1]

\[
S_n=\frac{10^n-1}{9}
\]
\hfill \textbf{A1}

\medskip

\textbf{(b)} Show that
$S_1+S_2+\cdots+S_n=\dfrac{10(10^n-1)-9n}{81}$. \hfill [4]

\[
S_1+S_2+\cdots+S_n
=\frac{10-1}{9}+\frac{10^2-1}{9}+\cdots+\frac{10^n-1}{9}
\]
\hfill \textbf{A1}

\[
=\frac{(10+10^2+\cdots+10^n)-n}{9}
\]
\hfill \textbf{M1}

\[
10+10^2+\cdots+10^n=\frac{10(10^n-1)}{10-1}, \qquad
(-1)+(-1)+\cdots+(-1)=-n
\]
\hfill \textbf{A1}

\[
S_1+S_2+\cdots+S_n
=\frac{\frac{10(10^n-1)}{9}-n}{9}
=\frac{10(10^n-1)-9n}{81}
\]
\hfill \textbf{A1 AG}

\newpage

%--- Problem 2: Geometric areas, mean, and median [8 marks] ---
%--- Source: AA SL 2025 May P1 TZ1 Q6 ---
\item[2.] [Maximum mark: 8]

\textbf{(a)(i)} Show that the area of $F_n$ is
$20\left(\dfrac{9}{4}\right)^{n-1}\text{ cm}^2$. \hfill [2]

The width and height each have common ratio $\dfrac{3}{2}$, so the areas
form a geometric sequence with first term $20$ and common ratio
\[
\left(\frac{3}{2}\right)^2=\frac{9}{4}.
\]
\hfill \textbf{M1 A1}

\[
A_n=20\left(\frac{9}{4}\right)^{n-1}
\]
\hfill \textbf{AG}

\medskip

\textbf{(a)(ii)} Find the mean area. \hfill [3]

\[
\sum_{n=1}^{10} A_n
=\frac{20\left(\left(\frac{9}{4}\right)^{10}-1\right)}
{\frac{9}{4}-1}
\]
\hfill \textbf{M1}

\[
=16\left(\left(\frac{9}{4}\right)^{10}-1\right)
\]
\hfill \textbf{A1}

Therefore the mean area is
\[
\frac{1}{10}\cdot 16\left(\left(\frac{9}{4}\right)^{10}-1\right)
=\frac{8}{5}\left(\left(\frac{9}{4}\right)^{10}-1\right)\text{ cm}^2.
\]
\hfill \textbf{A1}

\medskip

\textbf{(b)} Find the median area. \hfill [3]

The median is the mean of the $5$th and $6$th areas. \hfill \textbf{M1}

\[
\text{median}
=\frac{20\left(\frac{9}{4}\right)^4+
20\left(\frac{9}{4}\right)^5}{2}
\]
\hfill \textbf{A1}

\[
=\frac{20\left(\frac{9}{4}\right)^4\left(1+\frac{9}{4}\right)}{2}
=\frac{65}{2}\left(\frac{9}{4}\right)^4\text{ cm}^2.
\]
\hfill \textbf{A1}

\newpage

%--- Problem 3: Card pyramids and arithmetic sums [16 marks] ---
%--- Source: AA SL 2025 May P2 TZ1 Q8 ---
\item[3.] [Maximum mark: 16]

\textbf{(a)} Write down $t_3$. \hfill [1]

\[
t_3=15
\]
\hfill \textbf{A1}

\medskip

\textbf{(b)} Find $t_4$. \hfill [2]

\[
t_4=15+11
\]
\hfill \textbf{M1}

\[
t_4=26
\]
\hfill \textbf{A1}

\medskip

\textbf{(c)} Show that $t_n=\dfrac{n(3n+1)}{2}$. \hfill [3]

The row totals form the arithmetic sequence
\[
2+5+8+11+\cdots .
\]
\hfill \textbf{M1}

\[
t_n=\frac{n}{2}\left(2(2)+3(n-1)\right)
\]
\hfill \textbf{M1}

\[
t_n=\frac{n}{2}(3n+1)=\frac{n(3n+1)}{2}
\]
\hfill \textbf{A1 AG}

\medskip

\textbf{(d)} Find the maximum number of rows from $14$ full packs.
\hfill [3]

\[
\frac{n(3n+1)}{2}\leq 14(52)=728
\]
\hfill \textbf{M1}

Solving gives $n=21.8642\ldots$, or checking:
\[
t_{21}=672,\qquad t_{22}=737.
\]
\hfill \textbf{A1}

The maximum number of rows is $21$. \hfill \textbf{A1}

\medskip

\textbf{(e)} Find the minimum number of rows using full packs with no
cards left over. \hfill [2]

We need
\[
\frac{\frac{1}{2}n(3n+1)}{52}\in\mathbb{Z}.
\]
\hfill \textbf{M1}

The minimum number of rows is $13$. \hfill \textbf{A1}

\medskip

\textbf{(f)} Find the minimum number of cards for height more than
$2$ metres. \hfill [5]

The vertical height of one triangular row is
\[
88\sin 60^\circ=44\sqrt{3}\text{ mm}.
\]
\hfill \textbf{M1 A1}

We need
\[
44n\sqrt{3}>2000.
\]
\hfill \textbf{M1}

\[
n>26.2431\ldots,
\]
so the minimum number of rows is $27$. \hfill \textbf{A1}

\[
t_{27}=\frac{27(3(27)+1)}{2}=1107.
\]
\hfill \textbf{A1}

\bigskip
\hrule
\bigskip

%--- Problem 4: Simple and compound interest comparison [12 marks] ---
%--- Source: AA SL 2025 Nov P2 TZ3 Q7 ---
\item[4.] [Maximum mark: 12]

\textbf{(a)(i)} Monthly interest for Andy. \hfill [1]

\[
5000(0.00315)=15.75
\]
\hfill \textbf{A1}

\medskip

\textbf{(a)(ii)} Annual interest for Andy. \hfill [1]

\[
12(15.75)=189
\]
\hfill \textbf{A1}

\medskip

\textbf{(a)(iii)} Andy's amount after $n$ years. \hfill [1]

\[
5000+189n
\]
\hfill \textbf{A1}

\medskip

\textbf{(a)(iv)} Show Andy has $\$5945$ after $5$ years. \hfill [1]

\[
5000+189(5)=5945
\]
\hfill \textbf{A1 AG}

\medskip

\textbf{(b)(i)} Jess's account after $5$ years. \hfill [2]

\[
5000(1.03)^5
\]
\hfill \textbf{M1}

\[
=5796.37\ldots \approx \$5796
\]
\hfill \textbf{A1}

\medskip

\textbf{(b)(ii)} Jess's interest over $5$ years. \hfill [1]

\[
5796.37\ldots-5000=796.37\ldots \approx \$796
\]
\hfill \textbf{A1}

\medskip

\textbf{(c)(i)} Jess's amount after $n$ years. \hfill [2]

Use the compound-interest model with annual compounding. \hfill \textbf{M1}

\[
5000(1.03)^n
\]
\hfill \textbf{A1}

\medskip

\textbf{(c)(ii)} Smallest number of complete years. \hfill [3]

\[
5000(1.03)^n>5000+189n
\]
\hfill \textbf{M1}

Solving or checking values:
\[
5000(1.03)^{16}=8023.53\ldots<8024,
\qquad
5000(1.03)^{17}=8264.23\ldots>8213.
\]
\hfill \textbf{M1}

Therefore the smallest number of complete years is $17$. \hfill \textbf{A1}

\end{enumerate}
\end{document}
